
\section*{Ряд экспоненты и конечная формула Тейлора}
\begin{theorem}[Экспоненциальный ряд]
Для каждого действительного \(x\) выполнено \(e^x=\sum_{k=0}^\infty x^k/k!\).
\end{theorem}
\begin{proof}Повторное интегрирование по частям даёт конечную формулу
\[
e^x=\sum_{k=0}^N\frac{x^k}{k!}+\frac1{N!}\int_0^x e^t(x-t)^N\,dt.
\]
При \(N=0\) это Ньютона — Лейбница; один шаг интегрирования по частям извлекает следующее слагаемое, что доказывает формулу индукцией. Для любого знака \(x\) модуль остатка не больше \(e^{|x|}|x|^{N+1}/(N+1)!\). При \(x=0\) равенство непосредственно. При \(x\ne0\) последняя величина стремится к нулю: отношение соседних членов равно \(|x|/(N+2)\) и начиная с некоторого номера не больше \(1/2\). Значит, частичные суммы сходятся к \(e^x\).\end{proof}
Для произвольной функции тот же приём даёт формулу Тейлора с интегральным остатком, если производные до нужного порядка непрерывны на отрезке между \(0\) и \(x\). Здесь все производные экспоненты совпадают с ней; поэтому оценка остатка особенно проста.
